\documentclass[border=2pt]{standalone}
% Shared preamble for standalone figures in this directory.
% Compile with XeLaTeX from 讲义/figs, or via the figures target.
\usepackage{ctex}
\usepackage{amsmath}
\usepackage{amssymb}
\usepackage{bm}
\usepackage{mathtools}
\usepackage{tikz}
\renewcommand{\vec}[1]{\boldsymbol{#1}}
\newcommand{\cC}{\mathcal{C}}
\newcommand{\cD}{\mathcal{D}}
\newcommand{\R}{\mathbf{R}}
\usetikzlibrary{
    arrows.meta,
    calc,
    decorations.markings,
    angles,
    quotes,
    positioning,
    patterns,
    shapes.geometric
}

\begin{document}
\begin{tikzpicture}[
        >=Stealth,
        font=\small,
        box/.style={draw, align=center, inner sep=4.5pt},
        oval/.style={draw, ellipse, align=center, inner xsep=7pt, inner ysep=4pt},
        arr/.style={->, thick},
        both/.style={<->, thick}
    ]
    \node[oval] (dep) at (0,0) {线性相关性};
    \node[box] (comb) at (0.15,1.85) {线性组合\\（定义）};
    \node[box] (eq) at (-2.55,1.2) {线性方程组};
    \node[box] (part) at (-3.85,-0.05) {部分组\\线性相关性};
    \node[box, text width=6.5cm] (repr) at (-1.55,-2.45)
        {线性表示\\表示其他向量唯一？\\内部向量可以互相线性表示？\\存在后面的向量可以被前面的向量表示？};
    \node[box] (source) at (5.4,-1.0) {源泉定理};
    \node[box, text width=4.15cm] (pair) at (5.4,-2.65)
        {向量组 1 可以表示向量组 2\\二者秩的关系？};
    \node[oval] (max) at (-3.05,-4.85) {极大线性无关组};
    \node[oval] (rank) at (1.45,-4.85) {向量组的秩};
    \node[oval] (equiv) at (5.75,-4.85) {等价向量组};
    \node[box] (solve) at (-5.05,-6.55) {求解：初等变换};
    \node[oval] (basis) at (0.2,-7.05) {基和维数};
    \node[above] at (basis.north) {线性空间};
    \node[box] (nvec) at (-2.85,-8.65) {$n$ 维空间 $n$ 个\\线性无关向量};
    \node[box] (ext) at (0.2,-8.65) {基的扩充};
    \node[oval] (coord) at (3.05,-8.65) {向量的坐标};

    \draw[both] (comb) -- (dep);
    \draw[both] (eq.south east) -- (dep.north west);
    \draw[both] (part) -- (dep);
    \draw[both] (dep.south west) -- (repr.north);
    \draw[arr] (dep.east) -- (source.north);
    \draw[arr] (source) -- (pair);
    \draw[arr] (repr.210) -- node[left] {存在} (max.north);
    \draw[arr] (solve.north) -- (max.south west);
    \draw[arr] (max) -- (rank);
    \draw[arr] (rank.north east) -- (pair.210);
    \draw[arr] (pair.south) -- (equiv.north);
    \draw[arr] (max.south east) -- (basis.north west);
    \draw[arr] (rank.south east) -- (basis.north east);
    \draw[arr] (basis.south west) -- (nvec.north east);
    \draw[arr] (basis.south) -- (ext.north);
    \draw[arr] (basis.south east) -- (coord.north west);
\end{tikzpicture}
\end{document}
